% =============================================================================
%  DC³S: Telemetry-Reliability-Weighted Conformal Safety Shield for
%        Battery Dispatch under Degraded Mixed Telemetry
%  R1 Conference Submission Draft
% =============================================================================
\documentclass[11pt,a4paper,twocolumn]{article}

\usepackage[utf8]{inputenc}
\usepackage[T1]{fontenc}
\usepackage[margin=1in]{geometry}
\usepackage{booktabs}
\usepackage{longtable}
\usepackage{array}
\usepackage[hypertexnames=false]{hyperref}
\usepackage{amsmath,amssymb,amsthm}
\usepackage{graphicx}
\usepackage{enumitem}
\usepackage{xcolor}
\usepackage{algorithm}
\usepackage{algpseudocode}
\usepackage{multirow}
\usepackage{subcaption}
\usepackage{cleveref}

% ---------- Title ----------------------------------------------------------
\title{DC3S: Telemetry-Reliability-Weighted Conformal Safety Shield for Battery Dispatch under Degraded Mixed Telemetry}
\author{Pratik Niroula\\
        Minnesota State University, Mankato\\
        CIT (Major), Mathematics (Minor)\\
        \texttt{pratik.niroula@mnsu.edu}}
\date{}

% ---------- Macros ---------------------------------------------------------
\makeatletter
\newcommand{\DeclarePaperFigure}[2]{\expandafter\def\csname PaperFigurePath@#1\endcsname{#2}}
\newcommand{\DeclarePaperTable}[2]{\expandafter\def\csname PaperTablePath@#1\endcsname{#2}}
\newcommand{\PaperMetric}[2]{\@ifundefined{PaperMetric@#1}{#2}{\csname PaperMetric@#1\endcsname}}

\newcommand{\PaperFigure}[2]{%
  \begin{figure}[htbp]
  \centering
  \@ifundefined{PaperFigurePath@#1}{%
    \fbox{\parbox{0.9\linewidth}{\textbf{Missing figure token:} #1}}%
  }{%
    \IfFileExists{\csname PaperFigurePath@#1\endcsname}{%
      \includegraphics[width=0.9\linewidth]{\csname PaperFigurePath@#1\endcsname}%
    }{%
      \fbox{\parbox{0.9\linewidth}{\textbf{Missing figure file for token #1:} \csname PaperFigurePath@#1\endcsname}}%
    }%
  }%
  \caption{#2}
  \label{fig:#1}
  \end{figure}
}

\newcommand{\PaperTableCSV}[2]{%
  \@ifundefined{PaperTablePath@#1}{%
    \begin{table}[htbp]
    \centering
    \caption{#2}
    \fbox{\parbox{0.9\linewidth}{\textbf{Missing table token:} #1}}
    \label{tab:#1}
    \end{table}
  }{%
    \IfFileExists{\csname PaperTablePath@#1\endcsname}{%
      \InputIfFileExists{\csname PaperTablePath@#1\endcsname}{}{}%
    }{%
      \begin{table}[htbp]
      \centering
      \caption{#2}
      \fbox{\parbox{0.9\linewidth}{\textbf{Missing table file for token #1:} \csname PaperTablePath@#1\endcsname}}
      \label{tab:#1}
      \end{table}
    }%
  }%
}
\makeatother

% ---------- Locked metric values -------------------------------------------
\IfFileExists{assets/data/manifest_lookup.tex}{% Auto-generated by scripts/build_paper_table_tex.py
\DeclarePaperFigure{FIG01_ARCHITECTURE}{assets/figures/fig01_architecture.png}
\DeclarePaperFigure{FIG02_DC3S_STEP}{assets/figures/fig02_dc3s_step.png}
\DeclarePaperFigure{FIG03_TRUE_SOC_VIOLATION}{assets/figures/fig03_true_soc_violation_vs_dropout.png}
\DeclarePaperFigure{FIG04_TRUE_SOC_SEVERITY}{assets/figures/fig04_true_soc_severity_p95_vs_dropout.png}
\DeclarePaperFigure{FIG05_CQR_GROUP_COVERAGE}{assets/figures/fig05_cqr_group_coverage.png}
\DeclarePaperFigure{FIG06_TRANSFER_COVERAGE}{assets/figures/fig06_transfer_coverage.png}
\DeclarePaperFigure{FIG07_COST_SAFETY_FRONTIER}{assets/figures/fig07_cost_safety_frontier.png}
\DeclarePaperFigure{FIG08_RAC_SENSITIVITY_WIDTH}{assets/figures/fig08_rac_sensitivity_vs_width.png}
\DeclarePaperFigure{FIG09_REGION_DATASET_CARDS}{assets/figures/fig09_region_dataset_cards.png}
\DeclarePaperFigure{FIG10_CALIBRATION_TRADEOFF}{assets/figures/fig10_calibration_tradeoff.png}
\DeclarePaperFigure{FIG11_TRANSFER_GENERALIZATION}{assets/figures/fig11_transfer_generalization.png}
\DeclarePaperFigure{FIG12_GEOGRAPHIC_SCOPE}{assets/figures/fig12_geographic_scope.png}
\DeclarePaperFigure{FIG13_LOAD_RENEWABLE_PROFILES}{assets/figures/fig13_load_renewable_profiles.png}
\DeclarePaperFigure{FIG58_GAP_DISTRIBUTION}{assets/figures/fig58_gap_distribution.png}
\DeclarePaperFigure{FIG_UNIVERSAL_FRAMEWORK}{assets/figures/fig_universal_framework.png}
\DeclarePaperFigure{FIG14_FORECAST_VS_ACTUAL}{assets/figures/fig14_forecast_vs_actual.png}
\DeclarePaperFigure{FIG15_ROLLING_BACKTEST_RMSE}{assets/figures/fig15_rolling_backtest_rmse.png}
\DeclarePaperFigure{FIG16_ERROR_SEASONALITY}{assets/figures/fig16_error_seasonality.png}
\DeclarePaperFigure{FIG17_CONFORMAL_INTERVALS}{assets/figures/fig17_conformal_intervals.png}
\DeclarePaperFigure{FIG18_COVERAGE_VS_HORIZON}{assets/figures/fig18_coverage_vs_horizon.png}
\DeclarePaperFigure{FIG19_ANOMALY_TIMELINE}{assets/figures/fig19_anomaly_timeline.png}
\DeclarePaperFigure{FIG20_DISPATCH_COMPARISON}{assets/figures/fig20_dispatch_comparison.png}
\DeclarePaperFigure{FIG21_SOC_TRAJECTORY}{assets/figures/fig21_soc_trajectory.png}
\DeclarePaperFigure{FIG22_COST_CARBON_TRADEOFF}{assets/figures/fig22_cost_carbon_tradeoff.png}
\DeclarePaperFigure{FIG23_SAVINGS_SENSITIVITY}{assets/figures/fig23_savings_sensitivity.png}
\DeclarePaperFigure{FIG24_REGRET_PERTURBATION}{assets/figures/fig24_regret_perturbation.png}
\DeclarePaperFigure{FIG25_DATA_DRIFT}{assets/figures/fig25_data_drift.png}
\DeclarePaperFigure{FIG26_ABLATION_SENSITIVITY}{assets/figures/fig26_ablation_sensitivity.png}
\DeclarePaperFigure{FIG27_VIOLATION_VS_COST}{assets/figures/fig27_violation_vs_cost.png}
\DeclarePaperFigure{FIG28_DISTRIBUTIONAL_VS_CQR}{assets/figures/fig28_distributional_vs_cqr.png}
\DeclarePaperFigure{FIG29_CROSS_REGION_TRANSFER}{assets/figures/fig29_cross_region_transfer.png}
\DeclarePaperFigure{FIG30_VIOLATION_RATE}{assets/figures/fig30_violation_rate.png}
\DeclarePaperFigure{FIG31_VIOLATION_SEVERITY_P95}{assets/figures/fig31_violation_severity_p95.png}
\DeclarePaperFigure{FIG32_SHAP_SUMMARY_LOAD}{assets/figures/fig32_shap_summary_load.png}
\DeclarePaperFigure{FIG33_SHAP_SUMMARY_WIND}{assets/figures/fig33_shap_summary_wind.png}
\DeclarePaperFigure{FIG34_SHAP_SUMMARY_SOLAR}{assets/figures/fig34_shap_summary_solar.png}
\DeclarePaperFigure{FIG35_FORECAST_VS_ACTUAL_LOAD}{assets/figures/fig35_forecast_vs_actual_load.png}
\DeclarePaperFigure{FIG36_FORECAST_VS_ACTUAL_WIND}{assets/figures/fig36_forecast_vs_actual_wind.png}
\DeclarePaperFigure{FIG37_FORECAST_VS_ACTUAL_SOLAR}{assets/figures/fig37_forecast_vs_actual_solar.png}
\DeclarePaperFigure{FIG38_MODEL_COMPARISON}{assets/figures/fig38_model_comparison.png}
\DeclarePaperFigure{FIG39_DISPATCH_COMPARE}{assets/figures/fig39_dispatch_compare.png}
\DeclarePaperFigure{FIG40_ARBITRAGE_OPTIMIZATION}{assets/figures/fig40_arbitrage_optimization.png}
\DeclarePaperFigure{FIG41_SOC_TRAJECTORY_DETAIL}{assets/figures/fig41_soc_trajectory_detail.png}
\DeclarePaperFigure{FIG42_PREDICTION_INTERVALS_LOAD}{assets/figures/fig42_prediction_intervals_load.png}
\DeclarePaperFigure{FIG43_INTERVAL_WIDTH_BY_HORIZON}{assets/figures/fig43_interval_width_by_horizon.png}
\DeclarePaperFigure{FIG44_DATA_DRIFT_KS}{assets/figures/fig44_data_drift_ks.png}
\DeclarePaperFigure{FIG45_MODEL_DRIFT_METRIC}{assets/figures/fig45_model_drift_metric.png}
\DeclarePaperFigure{FIG46_ORIUS_PROGRAM}{assets/figures/fig46_orius_program_spine.png}
\DeclarePaperFigure{FIG47_ERROR_DISTRIBUTION}{assets/figures/fig47_error_distribution.png}
\DeclarePaperFigure{FIG48_SEASONALITY_ERROR_HEATMAP}{assets/figures/fig48_seasonality_error_heatmap.png}
\DeclarePaperFigure{FIG49_ROBUSTNESS_REGRET_BOXPLOT}{assets/figures/fig49_robustness_regret_boxplot.png}
\DeclarePaperFigure{FIG50_ROBUSTNESS_INFEASIBLE_RATE}{assets/figures/fig50_robustness_infeasible_rate.png}
\DeclarePaperFigure{FIG51_COVERAGE_BY_HORIZON}{assets/figures/fig51_coverage_by_horizon.png}
\DeclarePaperFigure{FIG52_COST_CARBON_DETAIL}{assets/figures/fig52_cost_carbon_detail.png}
\DeclarePaperFigure{FIG53_CASE_STUDY_DISPATCH}{assets/figures/fig53_case_study_dispatch.png}
\DeclarePaperFigure{FIG54_CASE_STUDY_WEEK}{assets/figures/fig54_case_study_week.png}
\DeclarePaperFigure{FIG55_MULTI_HORIZON_BACKTEST}{assets/figures/fig55_multi_horizon_backtest.png}
\DeclarePaperFigure{FIG56_USA_DATA_QUALITY}{assets/figures/fig56_usa_data_quality.png}
\DeclarePaperFigure{FIG57_GAP_HOURLY_PROFILES}{assets/figures/fig57_gap_hourly_profiles.png}
\DeclarePaperFigure{FIG_MULTI_DOMAIN_VALIDATION}{assets/figures/fig_multi_domain_validation.png}
\DeclarePaperTable{TBL01_MAIN_RESULTS}{assets/tables/generated/tbl01_main_results.tex}
\DeclarePaperTable{TBL02_ABLATIONS}{assets/tables/generated/tbl02_ablations.tex}
\DeclarePaperTable{TBL03_CQR_GROUP_COVERAGE}{assets/tables/generated/tbl03_cqr_group_coverage.tex}
\DeclarePaperTable{TBL04_TRANSFER_STRESS}{assets/tables/generated/tbl04_transfer_stress.tex}
\DeclarePaperTable{TBL05_DATASET_SUMMARY}{assets/tables/generated/tbl05_dataset_summary.tex}
\DeclarePaperTable{TBL06_HYPERPARAMS}{assets/tables/generated/tbl06_hyperparams.tex}
\DeclarePaperTable{TBL07_DATASET_CARDS}{assets/tables/generated/tbl07_dataset_cards.tex}
\DeclarePaperTable{TBL08_FORECAST_BASELINES}{assets/tables/generated/tbl08_forecast_baselines.tex}
\DeclarePaperTable{TBL09_CROSS_DOMAIN_OASG}{assets/tables/generated/cross_domain_oasg_table.tex}
\DeclarePaperTable{TBL10_DOMAIN_VALIDATION_STATUS}{assets/tables/generated/domain_validation_summary.tex}
}{}
\IfFileExists{assets/data/metrics_values.tex}{% Auto-generated by scripts/update_paper_metrics.py
}{}

% ---------- Theorem environments -------------------------------------------
\newtheorem{theorem}{Theorem}
\newtheorem{proposition}{Proposition}
\newtheorem{definition}{Definition}

% =============================================================================
\begin{document}
\maketitle

% ===== ABSTRACT ==============================================================
\begin{abstract}
Battery storage dispatch under degraded telemetry exposes a critical safety gap: when sensor readings are delayed, dropped, or corrupted, uncertainty estimates become miscalibrated and optimization solvers issue unsafe commands.  We introduce DC\textsuperscript{3}S (Detect--Calibrate--Constrain--Certify--Shield), a conformal safety shield that scores telemetry reliability at each control step, inflates prediction intervals via a Reliability-Adaptive Conformal Certification (RAC-Cert) law, projects candidate dispatch actions into the feasible safe set, and emits an auditable certificate.  Evaluated on a CPSBench harness covering four fault dimensions (dropout, delay-jitter, spikes, out-of-order) across Germany (OPSD, 17{,}377 rows) and three US balancing authorities (EIA-930 MISO/PJM/ERCOT, ${\sim}$13{,}500 rows each), DC\textsuperscript{3}S achieves \textbf{zero true-SOC violations} with a 2.8\% intervention rate while the deterministic baseline violates at 3.9\% with P95 severity of 333\,MWh.  Cost-side, DC\textsuperscript{3}S-wrapped dispatch yields \PaperMetric{DE_COST_SAVINGS_PCT}{7.11}\% cost savings, \PaperMetric{DE_CARBON_REDUCTION_PCT}{0.30}\% carbon reduction, and \PaperMetric{DE_PEAK_SHAVING_PCT}{6.13}\% peak shaving on DE; stochastic value of the stochastic solution is positive in both regions (DE VSS\,=\,\PaperMetric{DE_VSS}{2,708.61}, US VSS\,=\,\PaperMetric{US_VSS}{297,092.71}).  A nine-row ablation with Wilcoxon signed-rank gates isolates RAC-Cert inflation as the critical safety component.  All quantitative claims are locked to versioned run artifacts and validated by an automated claim-consistency checker.
\end{abstract}

\noindent\textbf{Keywords:} conformal prediction, battery dispatch, safety shield, telemetry reliability, robust optimization, cyber-physical systems, uncertainty quantification

% ===== 1. INTRODUCTION ======================================================
\section{Introduction}\label{sec:intro}

Grid-scale battery energy storage systems (BESS) are essential for integrating variable renewable generation, yet their dispatch depends on real-time telemetry---state-of-charge readings, power measurements, and grid-frequency signals---that is routinely degraded by sensor dropout, communication delay, spike injection, and packet reordering in cyber-physical power systems~\cite{dibaji2019systems,chandola2009anomaly}. When telemetry quality drops, the prediction intervals that inform optimization widen unpredictably, and dispatch commands computed from stale or corrupted state can push the battery into unsafe SOC regions, risking equipment damage and grid instability~\cite{chen2020review_bess}.

Existing approaches address parts of this pipeline in isolation.  Conformal prediction provides distribution-free coverage guarantees~\cite{vovk2005conformal,shafer2008tutorial}, and adaptive variants track distributional shift~\cite{gibbs2021adaptive,zaffran2022adaptive}, but neither incorporates telemetry-quality signals into the calibration step.  Robust and stochastic dispatch formulations protect against load/price uncertainty~\cite{bertsimas2011theory,birge2011stochastic}, but assume the uncertainty sets themselves are well-calibrated---an assumption violated under degraded telemetry.  Safety shields in control theory enforce state constraints~\cite{alshiekh2018safe,bastani2021safe}, but typically do not reason about measurement reliability or provide dispatch-level audit trails.

\textbf{This paper bridges these gaps} with DC\textsuperscript{3}S, a conformal safety shield that:
\begin{enumerate}[leftmargin=*,topsep=2pt,itemsep=1pt]
  \item \textbf{Detects} telemetry degradation via a composite reliability score $w_t$;
  \item \textbf{Calibrates} prediction intervals through RAC-Cert, inflating width proportional to degradation severity and drift risk;
  \item \textbf{Constrains} dispatch by solving a robust optimization problem over the inflated intervals;
  \item \textbf{Certifies} each dispatch action with a signed certificate recording the reliability state, interval widths, solver status, and constraint margins;
  \item \textbf{Shields} the battery by projecting any infeasible candidate action to the nearest safe command.
\end{enumerate}

\paragraph{Research questions.}
\begin{description}[leftmargin=1em,topsep=2pt,itemsep=1pt]
  \item[\textbf{RQ1}] Does RAC-Cert inflation eliminate true-SOC violations under telemetry faults without excessive cost penalty?
  \item[\textbf{RQ2}] How does DC\textsuperscript{3}S compare to deterministic, fixed-interval robust, and CVaR dispatch baselines on the cost--safety Pareto frontier?
  \item[\textbf{RQ3}] Which shield components (reliability scoring, interval inflation, action repair) are necessary---and which are sufficient---for safety?
  \item[\textbf{RQ4}] Do safety and cost outcomes transfer across regions (DE $\leftrightarrow$ US) and seasonal regimes?
\end{description}

\paragraph{Contributions.}
\begin{enumerate}[leftmargin=*,topsep=2pt,itemsep=1pt]
  \item \textbf{RAC-Cert}: a reliability-adaptive conformal certification law that inflates intervals as a function of telemetry quality, drift risk, and sensitivity, together with a supporting coverage-preservation proposition under monotone inflation (\Cref{thm:rac_coverage}).
  \item \textbf{CPSBench}: a fault-injection benchmark with truth-vs-observed SOC evaluation across four fault types and parametric severity sweeps.
  \item \textbf{Empirical evidence}: zero-violation safety with quantified cost--safety tradeoffs on DE and three US balancing authorities, including ablation and transfer studies.
  \item \textbf{Governance stack}: automated metric-locking and claim-validation tooling that prevents manuscript--artifact drift, reported as an explicit reproducibility contribution rather than background process~\cite{pineau2021improving}.
\end{enumerate}

% ===== 2. RELATED WORK ======================================================
\section{Related Work}\label{sec:related}

\paragraph{Conformal prediction for time series.}
Conformal prediction~\cite{vovk2005conformal,shafer2008tutorial} provides distribution-free prediction intervals with finite-sample coverage guarantees under exchangeability. Extensions to operational time-series settings include conformalized quantile regression (CQR)~\cite{romano2019conformalized}, adaptive conformal inference under shift~\cite{gibbs2021adaptive}, adaptive conformal prediction for time series~\cite{zaffran2022adaptive}, and conformal methods beyond exchangeability~\cite{barber2023conformal}. DC\textsuperscript{3}S inherits this foundation but departs from prior work by conditioning interval inflation on telemetry quality rather than on residual behavior alone.

\paragraph{Robust and stochastic battery dispatch.}
Robust optimization for energy dispatch~\cite{bertsimas2011theory,ben2009robust} hedges against worst-case uncertainty within prescribed sets. Stochastic programming~\cite{birge2011stochastic} optimizes expected cost over sampled scenarios, and CVaR-based formulations manage tail risk via Rockafellar--Uryasev linearization~\cite{rockafellar2000optimization}. In the BESS literature, Rosewater et al.~\cite{rosewater2020riskaverse} show that risk-averse model predictive control can improve storage dispatch robustness when model uncertainty is acknowledged explicitly. DC\textsuperscript{3}S is closer to this line of work than to generic storage scheduling papers, but targets a different failure mode: telemetry-induced miscalibration of uncertainty sets inside a dispatch loop that must remain safe under degraded measurements.

\paragraph{Safety shields and runtime monitors.}
Shield synthesis~\cite{alshiekh2018safe,konighofer2020shield} enforces safety by constraining the action space to precomputed safe regions. Model-predictive shielding~\cite{bastani2021safe} further tightens this connection between optimization and runtime safety. Simplex architectures~\cite{sha2001using} switch to a safe baseline controller upon detecting anomalies, while runtime verification monitors~\cite{leucker2009brief} check execution traces against safety properties. DC\textsuperscript{3}S combines these ideas at dispatch time: it monitors telemetry health, inflates the optimizer's admissible uncertainty set, repairs unsafe candidate actions, and emits a certificate that can be audited after execution.

\paragraph{Telemetry reliability in CPS.}
Telemetry degradation in industrial control has been studied in the context of CPS security~\cite{dibaji2019systems} and anomaly detection~\cite{chandola2009anomaly}, but these works stop short of turning telemetry reliability into a control input for storage dispatch. DC\textsuperscript{3}S uses the same threat and anomaly vocabulary---dropout, delay, spikes, stale or reordered observations---but pushes it into uncertainty calibration and action repair rather than treating it as a separate monitoring layer.

\paragraph{Reproducibility and release governance.}
Because this paper makes an explicit governance claim, the release process is positioned relative to the reproducibility standards discussed by Pineau et al.~\cite{pineau2021improving}. The goal is narrow: not to claim a new general theory of artifact governance, but to make run IDs, locked metrics, and manuscript validation part of the scientific object for this paper.

\paragraph{Positioning.}
\Cref{tab:positioning} summarizes the methodological positioning. DC\textsuperscript{3}S combines conformal calibration, telemetry-reliability weighting, dispatch-level action repair, and per-step certification in one framework. That combination, rather than any single ingredient in isolation, is the methodological point of departure from prior work.

\begin{table}[htbp]
\centering\small
\caption{Methodological positioning of DC\textsuperscript{3}S relative to prior work.}
\label{tab:positioning}
\begin{tabular}{lccccc}
\toprule
Method & \rotatebox{70}{Conformal} & \rotatebox{70}{Telemetry-aware} & \rotatebox{70}{Dispatch optim.} & \rotatebox{70}{Action repair} & \rotatebox{70}{Certificates} \\
\midrule
CQR / ACI        & \checkmark &            &            &            &            \\
Robust dispatch   &            &            & \checkmark &            &            \\
Safety shields    &            &            &            & \checkmark &            \\
Runtime monitors  &            &            &            &            & \checkmark \\
\textbf{DC\textsuperscript{3}S (ours)} & \checkmark & \checkmark & \checkmark & \checkmark & \checkmark \\
\bottomrule
\end{tabular}
\end{table}

\section{Background and Problem Formulation}\label{sec:background}

\subsection{System Architecture}\label{sec:system}

GridPulse is implemented as a closed forecast-to-dispatch stack with DC\textsuperscript{3}S inserted at the boundary between uncertainty estimation and control execution. Forecast models produce point and interval predictions for load, wind, and solar; the optimizer consumes those predictions to compute candidate battery actions; the shield then scores telemetry quality, widens uncertainty bounds when necessary, repairs unsafe actions, and records an auditable certificate. This architecture makes the safety argument operational rather than purely theoretical because the same control step produces both a dispatch command and the evidence needed to inspect it later.

\PaperFigure{FIG01_ARCHITECTURE}{GridPulse architecture with DC\textsuperscript{3}S at the dispatch boundary. The shield wraps the optimizer, inflating intervals and repairing actions before execution.}

This section plays two roles. First, it provides the background and preliminaries needed to read the formal problem statement and the later control sections. Second, it anchors the method in the actual repo-managed forecasting, optimization, monitoring, and reporting stack used to generate the evidence in this paper. DC\textsuperscript{3}S is therefore not evaluated as an isolated theorem or toy controller; it is evaluated inside an executable system whose artifacts remain tied to versioned outputs rather than drifting across manuscript revisions.

\subsection{Notation and Setting}\label{sec:problem}

Consider a discrete-time battery dispatch problem over horizon $T$ with time step $\Delta t$\,=\,1\,h.  At each step $t$, the controller observes telemetry $\mathbf{z}_t \in \mathbb{R}^d$ (load, wind, solar, price, SOC readings) and must select a dispatch action $a_t \in \mathcal{A}$ (charge/discharge power) subject to physical constraints.

\begin{definition}[Telemetry degradation]
A degradation operator $\mathcal{D}: \mathbb{R}^d \to \mathbb{R}^d$ maps clean telemetry $\mathbf{z}_t$ to observed telemetry $\tilde{\mathbf{z}}_t = \mathcal{D}(\mathbf{z}_t)$, where $\mathcal{D}$ may introduce dropout (zeroing), delay (time-shift), spikes (additive outliers), or reordering (permutation of recent samples).
\end{definition}

\begin{definition}[Truth vs.\ observed SOC]
Let $\text{SOC}_t^{\mathrm{true}}$ be the actual battery state governed by physics, and $\text{SOC}_t^{\mathrm{obs}}$ the value reported by (possibly degraded) sensors.  The safety evaluation is on $\text{SOC}_t^{\mathrm{true}}$; the controller only sees $\text{SOC}_t^{\mathrm{obs}}$.
\end{definition}

\subsection{Battery Dynamics}

The true SOC evolves as:
\begin{equation}\label{eq:soc_dynamics}
  \text{SOC}_{t+1}^{\mathrm{true}} = \text{SOC}_t^{\mathrm{true}}
    + \frac{\eta_c \cdot [a_t]^+ - [a_t]^- / \eta_d}{E_{\max}} \cdot \Delta t,
\end{equation}
where $\eta_c, \eta_d \in (0,1]$ are charge/discharge efficiencies, $E_{\max}$ is the rated energy capacity, and $[\cdot]^+, [\cdot]^-$ denote positive and negative parts.

The feasible action set enforces SOC bounds, power limits, and ramp constraints:
\begin{equation}\label{eq:constraints}
  \mathcal{A}_t = \bigl\{ a \in [-P_{\max}, P_{\max}] :
    \text{SOC}^{\min} \le \text{SOC}_{t+1}^{\mathrm{true}}(a) \le \text{SOC}^{\max},\;
    |a - a_{t-1}| \le R_{\max} \bigr\}.
\end{equation}

\subsection{Objective}

The dispatch objective balances cost, carbon, peak demand, and battery degradation:
\begin{equation}\label{eq:objective}
  \min_{a_{1:T}} \sum_{t=1}^{T}
    \underbrace{\lambda_c \cdot c_t \cdot a_t}_{\text{energy cost}}
    + \underbrace{\lambda_g \cdot g_t \cdot [a_t]^+}_{\text{carbon}}
    + \underbrace{\lambda_p \cdot \max(0, L_t + a_t - P_{\text{cap}})}_{\text{peak}}
    + \underbrace{\lambda_d \cdot |a_t|}_{\text{degradation}},
\end{equation}
where $c_t$ is the energy price, $g_t$ is the grid carbon intensity, $L_t$ is the net load, and $\lambda_c, \lambda_g, \lambda_p, \lambda_d$ are weighting parameters.

\subsection{Uncertainty Sets and the Calibration Gap}

Standard conformal prediction constructs intervals $\hat{C}_t(\alpha) = [\hat{q}_t^{\alpha/2},\, \hat{q}_t^{1-\alpha/2}]$ satisfying $\Pr(Y_t \in \hat{C}_t) \ge 1 - \alpha$ under exchangeability.  When $\tilde{\mathbf{z}}_t \ne \mathbf{z}_t$ due to degradation, the conformity scores shift, and the coverage guarantee degrades:
\begin{equation}\label{eq:coverage_gap}
  \Pr\bigl(Y_t \in \hat{C}_t(\alpha) \mid \mathcal{D} \text{ active}\bigr) < 1 - \alpha.
\end{equation}
This \emph{calibration gap} is the safety-critical failure mode that DC\textsuperscript{3}S is designed to close.

\subsection{Datasets}

We evaluate on two regions (\Cref{tab:TBL05_DATASET_SUMMARY}):
\begin{itemize}[leftmargin=*,topsep=2pt,itemsep=1pt]
  \item \textbf{DE}: Germany OPSD, 17{,}377 hourly rows (Oct 2018--Sep 2020), with load, wind, and solar signals.
  \item \textbf{US}: EIA-930 for MISO (13{,}638 rows), PJM (13{,}639 rows), and ERCOT (13{,}453 rows), Jul 2024--Jan 2026, same three signals.
\end{itemize}
All signals have 100\% non-null coverage.  The DE/US split provides cross-region transfer evaluation with distinct operating regimes.

\PaperTableCSV{TBL05_DATASET_SUMMARY}{Dataset summary: regions, date ranges, signal counts, and coverage.}

% ===== 4. THE DC3S METHOD ===================================================
\section{The \texorpdfstring{DC\textsuperscript{3}S}{DC3S} Method}\label{sec:method}

DC\textsuperscript{3}S operates on a narrow premise: if telemetry quality deteriorates, interval construction should become more conservative before optimization is allowed to act on those signals.  This section formalizes the five-step pipeline---reliability scoring, interval inflation, robust dispatch, action repair, and certificate emission---that translates degradation evidence into safe dispatch commands.

\subsection{Telemetry Reliability Scoring}\label{sec:reliability}

At each time step $t$, DC\textsuperscript{3}S computes a reliability score $w_t \in [0,1]$:
\begin{equation}\label{eq:reliability}
  w_t = \frac{1}{|\mathcal{S}|}\sum_{s \in \mathcal{S}}
    \bigl(1 - \mathbb{1}[\text{stale}_s(t)] \cdot \gamma_{\text{stale}}\bigr)
    \cdot \bigl(1 - \mathbb{1}[\text{spike}_s(t)] \cdot \gamma_{\text{spike}}\bigr)
    \cdot \bigl(1 - \mathbb{1}[\text{missing}_s(t)] \cdot \gamma_{\text{miss}}\bigr),
\end{equation}
where $\mathcal{S}$ is the sensor set, $\gamma_{\cdot} \in (0,1]$ are penalty weights, and the indicator functions detect staleness (no update within $\tau_{\text{stale}}$ seconds), spikes (exceeding $k\sigma$ from rolling median), and missing values.

\subsection{RAC-Cert Interval Inflation}\label{sec:rac}

Reliability-Adaptive Conformal Certification (RAC-Cert) inflates the conformal prediction interval as a function of reliability, drift risk, and model sensitivity:
\begin{equation}\label{eq:rac_inflation}
  \hat{C}_t^{\text{RAC}}(\alpha) = \hat{C}_t(\alpha) \cdot \bigl(1 + \underbrace{\kappa_r \cdot (1 - w_t)}_{\text{reliability inflation}} + \underbrace{\kappa_d \cdot d_t}_{\text{drift inflation}} + \underbrace{\kappa_s \cdot s_t}_{\text{sensitivity inflation}}\bigr),
\end{equation}
where:
\begin{itemize}[leftmargin=*,topsep=2pt,itemsep=1pt]
  \item $\kappa_r > 0$ controls the reliability-driven inflation rate;
  \item $d_t$ is a drift indicator (e.g., KS-test $p$-value below threshold); $\kappa_d > 0$ scales drift response;
  \item $s_t$ is a sensitivity probe measuring how much the forecast changes under input perturbation; $\kappa_s > 0$ scales sensitivity response.
\end{itemize}

With shrinkage toward the base width at rate $\tau \in (0,1)$ when reliability recovers:
\begin{equation}\label{eq:rac_shrinkage}
  \hat{w}_t^{\text{RAC}} = \tau \cdot \hat{w}_{t-1}^{\text{RAC}} + (1 - \tau) \cdot \hat{w}_t^{\text{raw}},
\end{equation}
preventing oscillation between inflated and nominal widths.

\begin{proposition}[RAC-Cert marginal coverage under monotone inflation]\label{thm:rac_coverage}
If the base conformal predictor satisfies $\Pr(Y_t \in \hat{C}_t(\alpha)) \ge 1 - \alpha$ under the clean data distribution, and RAC-Cert only \emph{inflates} intervals (the multiplicative factor in \Cref{eq:rac_inflation} is $\ge 1$), then the RAC-inflated interval preserves marginal coverage:
\[
  \Pr(Y_t \in \hat{C}_t^{\mathrm{RAC}}(\alpha)) \ge 1 - \alpha.
\]
\end{proposition}
\begin{proof}
Since $\hat{C}_t(\alpha) \subseteq \hat{C}_t^{\mathrm{RAC}}(\alpha)$ by construction (inflation factor $\ge 1$), we have $\{Y_t \in \hat{C}_t(\alpha)\} \subseteq \{Y_t \in \hat{C}_t^{\mathrm{RAC}}(\alpha)\}$, hence $\Pr(Y_t \in \hat{C}_t^{\mathrm{RAC}}(\alpha)) \ge \Pr(Y_t \in \hat{C}_t(\alpha)) \ge 1 - \alpha$.
\end{proof}

\PaperFigure{FIG08_RAC_SENSITIVITY_WIDTH}{RAC-Cert sensitivity signal versus interval width expansion across controllers.  Higher degradation triggers larger inflation.}

After RAC-Cert inflation, DC\textsuperscript{3}S passes the widened uncertainty set to the optimizer, verifies feasibility, repairs unsafe candidate actions when necessary, and records the evidence for that decision.  This stage is where conformal uncertainty becomes a control object.

\subsection{Robust Dispatch with Inflated Intervals}\label{sec:dispatch}

Given the RAC-inflated intervals $[\ell_t^{\text{RAC}}, u_t^{\text{RAC}}]$ for load, wind, and solar, DC\textsuperscript{3}S solves a two-scenario robust dispatch:
\begin{equation}\label{eq:robust_dispatch}
  \min_{a_{1:T}} \max_{Y \in \{\ell^{\text{RAC}}, u^{\text{RAC}}\}}
    \sum_{t=1}^{T} J(a_t, Y_t),
\end{equation}
where $J(\cdot)$ is the multi-objective cost from \Cref{eq:objective}.  This is implemented via an epigraph reformulation solvable by LP/MIP solvers.

\paragraph{CVaR baseline.}
For comparison, we also implement the CVaR formulation via Rockafellar--Uryasev linearization:
\begin{equation}\label{eq:cvar}
  \min_{a,\eta} \; \eta + \frac{1}{(1-\beta)S}\sum_{s=1}^{S} u_s,
  \quad u_s \ge J_s(a) - \eta,\; u_s \ge 0,
\end{equation}
where $\beta$ is the CVaR confidence level and $S$ is the scenario count.

\subsection{Action Repair and Projection}\label{sec:repair}

If the dispatch candidate $a_t^*$ violates any constraint in $\mathcal{A}_t$ (\Cref{eq:constraints}), DC\textsuperscript{3}S projects to the nearest feasible action:
\begin{equation}\label{eq:projection}
  a_t^{\text{safe}} = \arg\min_{a \in \mathcal{A}_t} \|a - a_t^*\|_2.
\end{equation}
For box constraints, this reduces to simple clamping.  For coupled constraints (ramp + SOC), a single-step QP is solved.

\subsection{Certificate Emission}\label{sec:certificate}

Each dispatch step emits a certificate $\mathcal{C}_t$ recording:
\begin{equation}\label{eq:certificate}
  \mathcal{C}_t = \bigl(t,\; w_t,\; \hat{C}_t^{\text{RAC}},\; a_t^*,\; a_t^{\text{safe}},\;
    \text{solver\_status},\; \text{SOC}_t^{\text{obs}},\; \Delta_{\text{margin}}\bigr),
\end{equation}
where $\Delta_{\text{margin}}$ records the distance to the nearest constraint boundary.  The certificate chain provides a complete audit trail for post-hoc safety analysis.

\PaperFigure{FIG02_DC3S_STEP}{DC\textsuperscript{3}S runtime step: observe $\to$ score $\to$ inflate $\to$ solve $\to$ repair $\to$ certify.}

\subsection{Full Algorithm}

\begin{algorithm}[htbp]
\caption{DC\textsuperscript{3}S: One dispatch step}
\label{alg:dc3s}
\begin{algorithmic}[1]
\Require Telemetry $\tilde{\mathbf{z}}_t$, base conformal predictor $\hat{C}_t(\alpha)$, optimizer $\mathcal{O}$
\Ensure Safe dispatch action $a_t^{\text{safe}}$, certificate $\mathcal{C}_t$
\State $w_t \gets \textsc{ReliabilityScore}(\tilde{\mathbf{z}}_t)$ \Comment{\Cref{eq:reliability}}
\State $d_t \gets \textsc{DriftDetect}(\tilde{\mathbf{z}}_t)$ \Comment{KS-test on residuals}
\State $s_t \gets \textsc{SensitivityProbe}(\tilde{\mathbf{z}}_t)$ \Comment{perturbation response}
\State $\hat{C}_t^{\text{RAC}} \gets \hat{C}_t(\alpha) \cdot (1 + \kappa_r(1-w_t) + \kappa_d d_t + \kappa_s s_t)$ \Comment{\Cref{eq:rac_inflation}}
\State $a_t^* \gets \mathcal{O}(\hat{C}_t^{\text{RAC}})$ \Comment{solve robust dispatch, \Cref{eq:robust_dispatch}}
\State $a_t^{\text{safe}} \gets \textsc{Project}(a_t^*, \mathcal{A}_t)$ \Comment{\Cref{eq:projection}}
\State $\mathcal{C}_t \gets \textsc{EmitCert}(t, w_t, \hat{C}_t^{\text{RAC}}, a_t^*, a_t^{\text{safe}})$ \Comment{\Cref{eq:certificate}}
\State \Return $a_t^{\text{safe}}, \mathcal{C}_t$
\end{algorithmic}
\end{algorithm}

% ===== 5. EXPERIMENTAL SETUP ================================================
\section{Experimental Setup}\label{sec:experiments}

\subsection{Forecasting Layer}

Three model families are trained on each region:
\begin{itemize}[leftmargin=*,topsep=2pt,itemsep=1pt]
  \item \textbf{GBM}: LightGBM gradient-boosted trees with lag, rolling, and calendar features.
  \item \textbf{LSTM}: two-layer LSTM with 128 hidden units, 168-hour lookback.
  \item \textbf{TCN}: temporal convolutional network with 3 residual blocks.
\end{itemize}
Each model is trained on three targets (load, wind, solar) with 24-hour forecast horizon, 10-fold time-aware cross-validation, and 24-hour gap between folds.

\subsection{Conformal Calibration}

Conformalized quantile regression (CQR) wraps each forecast model:
\begin{enumerate}[leftmargin=*,topsep=2pt,itemsep=1pt]
  \item Train quantile models at levels $\alpha/2$ and $1-\alpha/2$ (default $\alpha=0.10$).
  \item Compute conformity scores on calibration fold: $s_i = \max(\hat{q}_i^{\alpha/2} - Y_i,\; Y_i - \hat{q}_i^{1-\alpha/2})$.
  \item Set conformal quantile $\hat{q} = \text{Quantile}_{1-\alpha}(\{s_i\})$.
  \item Prediction interval: $[L_i - \hat{q},\; U_i + \hat{q}]$.
\end{enumerate}
Regime-aware CQR further stratifies calibration into low/mid/high volatility bins using a rolling standard deviation window.

\PaperTableCSV{TBL03_CQR_GROUP_COVERAGE}{Regime-aware CQR coverage (PICP@90) and mean interval width by target and volatility group.}

\subsection{CPSBench: Fault Injection Harness}

CPSBench evaluates controllers under four telemetry fault types:
\begin{enumerate}[leftmargin=*,topsep=2pt,itemsep=1pt]
  \item \textbf{Dropout}: random zeroing of sensor readings at rates 10--50\%.
  \item \textbf{Delay-jitter}: time-shifting readings by 1--5 steps with random perturbation.
  \item \textbf{Spikes}: additive outliers at $3\sigma$--$10\sigma$ from the signal mean.
  \item \textbf{Out-of-order}: permuting the last $k \in \{2,\ldots,8\}$ readings.
\end{enumerate}
Each fault type is swept across 10 severity levels.  All evaluation is on \emph{true} SOC (ground-truth battery state), not the observed SOC---this is the critical design choice that prevents controllers from ``hiding'' safety violations behind sensor errors.

\subsection{Baselines and Controllers}

Five controllers are compared (\Cref{tab:TBL01_MAIN_RESULTS}):
\begin{enumerate}[leftmargin=*,topsep=2pt,itemsep=1pt]
  \item \textbf{Deterministic LP}: point-forecast dispatch without uncertainty.
  \item \textbf{Robust fixed-interval}: robust dispatch with static $\pm 2\sigma$ intervals.
  \item \textbf{CVaR interval}: CVaR dispatch with scenario sampling from conformal intervals.
  \item \textbf{DC\textsuperscript{3}S wrapped}: DC\textsuperscript{3}S with RAC-Cert inflation around the robust optimizer.
  \item \textbf{DC\textsuperscript{3}S FTIT}: full-time-in-trade variant with aggressive intervention.
\end{enumerate}

\subsection{Evaluation Metrics}

\begin{itemize}[leftmargin=*,topsep=2pt,itemsep=1pt]
  \item \textbf{Safety}: true-SOC violation rate (fraction of steps where $\text{SOC}^{\mathrm{true}} \notin [\text{SOC}^{\min}, \text{SOC}^{\max}]$) and P95 violation severity (MWh).
  \item \textbf{Cost}: cost delta \% relative to no-storage baseline.
  \item \textbf{Intervention rate}: fraction of steps where DC\textsuperscript{3}S modifies the optimizer's candidate action.
  \item \textbf{Stochastic value}: VSS (value of stochastic solution) and EVPI (expected value of perfect information).
  \item \textbf{UQ quality}: PICP, pinball loss, Winkler score per the UQ metric contract.
\end{itemize}

\subsection{Statistical Protocol}

All randomized experiments use 10 seeds.  Safety comparisons use Wilcoxon signed-rank tests with $p < 0.01$ threshold.  Ablations report 95\% bootstrap CIs and apply two-gate verification: a relative-reduction gate ($\ge 10\%$) and a significance gate ($p < 0.01$).

% ===== 6. RESULTS ============================================================
\section{Results}\label{sec:results}

\subsection{Main Result: Safety under Telemetry Faults}\label{sec:main_result}

\Cref{tab:TBL01_MAIN_RESULTS} presents the core safety--cost comparison across all five controllers under the aggregate fault sweep.

\PaperTableCSV{TBL01_MAIN_RESULTS}{Main controller comparison under aggregate telemetry faults.  DC\textsuperscript{3}S (wrapped and FTIT) achieve zero true-SOC violations.  CI: 95\% bootstrap over 10 seeds.}

\textbf{Key findings:}
\begin{itemize}[leftmargin=*,topsep=2pt,itemsep=1pt]
  \item \textbf{DC\textsuperscript{3}S achieves zero violations}: both DC\textsuperscript{3}S variants maintain 0.000 true-SOC violation rate across all fault types and severity levels, while the deterministic baseline violates at 3.9\% (CI: [3.0\%, 4.8\%]) with P95 severity of 333\,MWh.
  \item \textbf{Low intervention cost}: DC\textsuperscript{3}S FTIT intervenes at only 2.8\% of steps (CI: [2.2\%, 3.4\%]), demonstrating that safety does not require constant override.
  \item \textbf{CVaR is unsafe}: the CVaR baseline violates at 25.6\% with severity 0.28\,MWh---lower severity but far more frequent than acceptable.
  \item \textbf{Cost neutrality}: DC\textsuperscript{3}S wrapped shows a $-1.9$\% cost delta (modest cost increase), while FTIT achieves $+3.0$\% (cost \emph{improvement}) due to more conservative scheduling that avoids peak-hour penalties.
\end{itemize}

\PaperFigure{FIG03_TRUE_SOC_VIOLATION}{True-SOC violation rate versus fault severity.  DC\textsuperscript{3}S maintains zero violations across all severity levels while baselines deteriorate monotonically.}

\PaperFigure{FIG04_TRUE_SOC_SEVERITY}{P95 true-SOC violation severity under fault sweep.  The deterministic baseline's severity grows superlinearly with fault severity; DC\textsuperscript{3}S remains at zero.}

\subsection{Cost--Safety Frontier}\label{sec:frontier}

\Cref{fig:FIG07_COST_SAFETY_FRONTIER} visualizes the Pareto frontier.  DC\textsuperscript{3}S occupies the dominant position: zero violations with competitive cost.  The deterministic baseline sits in the unsafe region, and CVaR achieves moderate cost improvement but at unacceptable violation rates.

\PaperFigure{FIG07_COST_SAFETY_FRONTIER}{Cost--safety Pareto frontier.  DC\textsuperscript{3}S dominates the zero-violation, low-cost region.  Dashed line: acceptable violation threshold.}

\subsection{Decision Impact: DE and US}\label{sec:impact}

Under locked artifact evaluation:
\begin{itemize}[leftmargin=*,topsep=2pt,itemsep=1pt]
  \item \textbf{DE}: \PaperMetric{DE_COST_SAVINGS_PCT}{7.11}\% cost savings, \PaperMetric{DE_CARBON_REDUCTION_PCT}{0.30}\% carbon reduction, \PaperMetric{DE_PEAK_SHAVING_PCT}{6.13}\% peak shaving.
  \item \textbf{US}: \PaperMetric{US_COST_SAVINGS_PCT}{0.11}\% cost savings, \PaperMetric{US_CARBON_REDUCTION_PCT}{0.13}\% carbon reduction, \PaperMetric{US_PEAK_SHAVING_PCT}{0.00}\% peak shaving.
\end{itemize}

The magnitude difference reflects operating-regime heterogeneity: DE's time-of-use tariff structure creates larger arbitrage opportunities than US wholesale markets in the evaluation window.  This supports the claim that forecast-to-dispatch systems should be evaluated as region-specific decision systems, not by aggregate forecast-score transferability alone.

\textbf{Stochastic value:}
\begin{itemize}[leftmargin=*,topsep=2pt,itemsep=1pt]
  \item DE: VSS\,=\,\PaperMetric{DE_VSS}{2,708.61}, EVPI\textsubscript{robust}\,=\,\PaperMetric{DE_EVPI_ROBUST}{2.32}, EVPI\textsubscript{det}\,=\,\PaperMetric{DE_EVPI_DETERMINISTIC}{-30.40}.
  \item US: VSS\,=\,\PaperMetric{US_VSS}{297,092.71}, EVPI\textsubscript{robust}\,=\,\PaperMetric{US_EVPI_ROBUST}{10,279,851.74}, EVPI\textsubscript{det}\,=\,\PaperMetric{US_EVPI_DETERMINISTIC}{24,915,503.93}.
\end{itemize}
Positive VSS in both regions confirms that uncertainty-aware dispatch provides measurable value over deterministic dispatch, answering RQ2.

\subsection{Ablation Study: What Matters for Safety}\label{sec:ablation}

\Cref{tab:TBL02_ABLATIONS} reports the nine-row ablation.  Each row removes one DC\textsuperscript{3}S component or replaces it with a simpler alternative.

\PaperTableCSV{TBL02_ABLATIONS}{Ablation study with strict two-gate verification (relative reduction $\ge$ 10\% \emph{and} Wilcoxon $p < 0.01$).}

\textbf{Key ablation findings:}
\begin{enumerate}[leftmargin=*,topsep=2pt,itemsep=1pt]
  \item \textbf{RAC-Cert is necessary}: removing reliability-adaptive inflation causes violations to reappear, failing the relative-reduction gate.
  \item \textbf{Action repair is necessary}: without projection, the optimizer occasionally produces infeasible commands that violate SOC bounds.
  \item \textbf{Drift detection is complementary}: removing drift inflation degrades severity P95 but does not cause violations at low fault levels---it becomes critical at high severity.
  \item \textbf{Fixed-interval robust is insufficient}: static intervals do not adapt to degradation level, leading to either excessive cost (wide intervals) or safety violations (narrow intervals).
\end{enumerate}
This answers RQ3: reliability scoring \emph{and} action repair are both necessary; drift detection is necessary for robustness under high-severity faults.

\subsection{Uncertainty Calibration Quality}\label{sec:calibration}

\Cref{tab:TBL03_CQR_GROUP_COVERAGE} reports regime-stratified CQR coverage.  Key observations:
\begin{itemize}[leftmargin=*,topsep=2pt,itemsep=1pt]
  \item Wind high-group achieves 90.0\% PICP (nominal target), but solar high-group drops to 73.8\%---an honest negative finding.
  \item Load low-group (85.2\%) and solar low-group (80.2\%) are below nominal, motivating the RAC-Cert inflation that compensates during dispatch.
  \item Mean interval widths range from 719--1{,}244\,MW, reflecting genuine forecast difficulty variation across targets and regimes.
\end{itemize}

\PaperFigure{FIG05_CQR_GROUP_COVERAGE}{Regime-aware CQR group coverage.  Below-nominal coverage in solar high-group is compensated by RAC-Cert inflation at dispatch time.}

\PaperFigure{FIG10_CALIBRATION_TRADEOFF}{Calibration tradeoff: coverage versus interval width across all target-regime combinations.}

\subsection{Transfer and Generalization}\label{sec:transfer}

Four transfer cases evaluate cross-region and temporal robustness:
\begin{enumerate}[leftmargin=*,topsep=2pt,itemsep=1pt]
  \item DE $\to$ US (no retrain): models trained on DE applied to US data.
  \item US $\to$ DE (no retrain): models trained on US applied to DE data.
  \item DE seasonal shift: evaluation on a held-out season.
  \item US seasonal shift: evaluation on a held-out season.
\end{enumerate}

\PaperTableCSV{TBL04_TRANSFER_STRESS}{Transfer stress outcomes: coverage, violation rate, severity, and cost delta across all four cases.}

\PaperFigure{FIG06_TRANSFER_COVERAGE}{Transfer stress coverage summary.  DC\textsuperscript{3}S maintains safety across transfer cases, though cost optimality degrades without retraining.}

\subsection{Runtime and Intervention Behavior}\label{sec:runtime}

DC\textsuperscript{3}S's computational overhead is modest: the reliability scoring, inflation calculation, and projection add $<$5\% to the solver wall-clock time. The intervention rate (fraction of steps where the safe action differs from the optimizer's candidate) is 2.8\% for DC\textsuperscript{3}S FTIT---meaning the optimizer's output is safe in $>$97\% of steps, and the shield activates only when genuinely needed.

% ===== 7. GOVERNANCE ========================================================
\section{Governance and Reproducibility}\label{sec:governance}

The paper's quantitative claims are locked to versioned artifacts rather than free-text manuscript statements. The locked DE and US impact summaries, the canonical stochastic run IDs, and the dashboard profile manifest are treated as publication interfaces. A claim is therefore valid only if it can be traced to a source artifact, a deterministic extraction rule, and a validator pass. This is the narrow sense in which governance is claimed as a contribution: reproducibility is embedded into the release procedure for this paper, consistent with the community push for stronger ML reproducibility norms~\cite{pineau2021improving}.

\paragraph{Run-scoped evidence.}
The release process binds the manuscript to dataset-scoped report files, pinned run IDs, and a claim matrix. This prevents a common failure mode in long-running applied ML projects: a paper draft continues to circulate after the underlying evaluation artifacts have changed. In this draft, DE and US stochastic results are always reported with their run IDs, and region-level impact percentages are always read from the locked artifact family rather than recomputed manually.

\paragraph{What is and is not claimed.}
The governance contribution is operational, not philosophical. We do not claim a new general reproducibility framework for ML. We claim that for this particular CPS paper, metric locking, run provenance, and non-mutating consistency checks are necessary to keep the method story and the evidence story aligned.

\paragraph{Repository artifact traceability.}
\Cref{tab:repo_traceability} makes the repo-to-paper mapping explicit. It mirrors the paper structure used by this draft: the method section is implemented by train/eval code, experiments consume configs, seeds, and preprocessed data, results are derived from logs and checkpoints, and the paper figures and tables are generated by plotting and publication scripts.

\begin{table}[htbp]
\centering\small
\caption{Repository artifact traceability for the paper pipeline.}
\label{tab:repo_traceability}
\begin{tabular}{p{3.1cm}p{4.2cm}p{5.8cm}}
\toprule
Repo artifact family & Primary role in the paper & Main paper touchpoint \\
\midrule
README + setup docs & Environment bootstrap, run commands, and operator setup & Governance and reproducibility; appendix checklist \\
Train/eval code & Implements forecasting, uncertainty calibration, dispatch, and CPSBench evaluation & Method and experimental setup \\
Configs + seeds & Fixes dataset scope, model settings, calibration settings, and randomized runs & Experimental setup and appendix hyperparameters \\
Data + preprocessing & Defines dataset construction, split logic, and feature generation & Background/preliminaries and experimental setup \\
Logs + checkpoints & Preserve run IDs, solver outputs, metrics, and locked evidence values & Results and governance/reproducibility \\
Plotting + publication scripts & Build figures, tables, and paper-facing assets from locked artifacts & Results and appendices \\
\bottomrule
\end{tabular}
\end{table}

% ===== 8. DISCUSSION =======================================================
\section{Discussion and Threats to Validity}\label{sec:discussion}

\subsection{Implications for CPS Safety}

DC\textsuperscript{3}S demonstrates that \emph{telemetry reliability should be a first-class input to the dispatch optimization loop}, not merely a pre-processing concern.  The zero-violation result is achieved not by being maximally conservative (which would sacrifice cost), but by \emph{adapting} conservatism to measurement quality.  This principle generalizes beyond battery dispatch to any CPS where control decisions depend on degraded sensor inputs: HVAC systems, water treatment, autonomous vehicles under sensor failure.

\subsection{Negative and Neutral Findings}

We report several findings that qualify our claims:
\begin{enumerate}[leftmargin=*,topsep=2pt,itemsep=1pt]
  \item \textbf{US peak shaving is 0.00\%} in locked evidence---the US wholesale market structure in the evaluation window did not create peak-shaving opportunities.
  \item \textbf{Solar high-regime PICP is 73.8\%}, below the 90\% nominal target.  RAC-Cert compensates during dispatch, but better base calibration would reduce unnecessary intervention.
  \item \textbf{Deep models (LSTM, TCN) do not outperform GBM} in locked dashboard snapshots---GBM dominates across all targets in both regions.
  \item \textbf{Cross-region impact magnitudes differ substantially}: a strategy yielding 7.11\% cost savings in DE yields only 0.11\% in US, underscoring regime sensitivity.
\end{enumerate}

\subsection{Threats to Validity and Limitations}\label{sec:limitations}

\begin{enumerate}[leftmargin=*,topsep=2pt,itemsep=1pt]
  \item \textbf{Dataset scope}: evidence is limited to DE (2018--2020) and US (2024--2026) windows.  Generalization to other markets, tariff structures, or regulatory regimes requires additional evaluation.
  \item \textbf{Simulation only}: all results are from software simulation with replayed telemetry.  No hardware-in-the-loop or field deployment validation is claimed.
  \item \textbf{Proxy assumptions}: cost and carbon signals are derived from publicly available data and may not fully represent all market settlement mechanisms.
  \item \textbf{Stochastic sensitivity}: VSS and EVPI magnitudes are sensitive to scenario construction parameters and should not be interpreted as universal constants.
  \item \textbf{Supporting proposition}: \Cref{thm:rac_coverage} provides marginal coverage preservation under inflation; it does not guarantee \emph{conditional} coverage or optimal interval width.
\end{enumerate}

% ===== 9. CONCLUSION =======================================================
\section{Conclusion}\label{sec:conclusion}

We presented DC\textsuperscript{3}S, a conformal safety shield that integrates telemetry-reliability scoring, adaptive interval inflation (RAC-Cert), robust dispatch, action repair, and per-step certification into a single online loop for battery dispatch under degraded telemetry.  Across four fault types and parametric severity sweeps on DE and US data:

\begin{itemize}[leftmargin=*,topsep=2pt,itemsep=1pt]
  \item DC\textsuperscript{3}S achieves \textbf{zero true-SOC violations} (vs.\ 3.9\% for deterministic and 25.6\% for CVaR baselines) with only 2.8\% intervention rate.
  \item Cost impact: \PaperMetric{DE_COST_SAVINGS_PCT}{7.11}\% cost savings (DE), \PaperMetric{US_COST_SAVINGS_PCT}{0.11}\% (US); positive VSS in both regions.
  \item Ablation confirms RAC-Cert inflation and action repair are both necessary for zero-violation safety.
\end{itemize}

The practical message: trustworthy CPS dispatch requires telemetry-aware uncertainty calibration, not just robust optimization over assumed-correct uncertainty sets.  DC\textsuperscript{3}S makes telemetry quality a \emph{control input} rather than a hidden failure mode.

\paragraph{Future work.}
\begin{enumerate}[leftmargin=*,topsep=2pt,itemsep=1pt]
  \item Extend to hardware-in-the-loop validation with physical inverter/BMS integration.
  \item Improve base calibration on under-covered targets (solar high-regime) to reduce unnecessary inflation.
  \item Expand evaluation to additional regions, tariff structures, and longer time windows.
  \item Investigate conditional coverage guarantees under non-stationary degradation.
\end{enumerate}

% ===== REFERENCES ============================================================
\bibliographystyle{plain}

\begin{thebibliography}{99}

\bibitem{vovk2005conformal}
V.~Vovk, A.~Gammerman, and G.~Shafer.
\newblock {\em Algorithmic Learning in a Random World}.
\newblock Springer, 2005.

\bibitem{shafer2008tutorial}
G.~Shafer and V.~Vovk.
\newblock A tutorial on conformal prediction.
\newblock {\em Journal of Machine Learning Research}, 9:371--421, 2008.

\bibitem{gibbs2021adaptive}
I.~Gibbs and E.~Cand\`{e}s.
\newblock Adaptive conformal inference under distribution shift.
\newblock {\em Advances in Neural Information Processing Systems}, 34, 2021.

\bibitem{zaffran2022adaptive}
M.~Zaffran, O.~Féron, Y.~Goude, J.~Josse, and A.~Dieuleveut.
\newblock Adaptive conformal predictions for time series.
\newblock {\em International Conference on Machine Learning}, 2022.

\bibitem{romano2019conformalized}
Y.~Romano, E.~Patterson, and E.~Cand\`{e}s.
\newblock Conformalized quantile regression.
\newblock {\em Advances in Neural Information Processing Systems}, 32, 2019.

\bibitem{barber2023conformal}
R.~F.~Barber, E.~J.~Cand\`{e}s, A.~Ramdas, and R.~J.~Tibshirani.
\newblock Conformal prediction beyond exchangeability.
\newblock {\em Annals of Statistics}, 51(2):816--845, 2023.

\bibitem{bertsimas2011theory}
D.~Bertsimas, D.~B.~Brown, and C.~Caramanis.
\newblock Theory and applications of robust optimization.
\newblock {\em SIAM Review}, 53(3):464--501, 2011.

\bibitem{ben2009robust}
A.~Ben-Tal, L.~El~Ghaoui, and A.~Nemirovski.
\newblock {\em Robust Optimization}.
\newblock Princeton University Press, 2009.

\bibitem{birge2011stochastic}
J.~R.~Birge and F.~Louveaux.
\newblock {\em Introduction to Stochastic Programming}.
\newblock Springer, 2nd edition, 2011.

\bibitem{rockafellar2000optimization}
R.~T.~Rockafellar and S.~Uryasev.
\newblock Optimization of conditional value-at-risk.
\newblock {\em Journal of Risk}, 2(3):21--42, 2000.

\bibitem{rosewater2020riskaverse}
D.~Rosewater, R.~Baldick, and S.~Santoso.
\newblock Risk-averse model predictive control design for battery energy storage systems.
\newblock {\em IEEE Transactions on Smart Grid}, 11(3):2014--2022, 2020.

\bibitem{alshiekh2018safe}
M.~Alshiekh, R.~Bloem, R.~Ehlers, B.~K\"onighofer, S.~Niekum, and U.~Topcu.
\newblock Safe reinforcement learning via shielding.
\newblock {\em AAAI Conference on Artificial Intelligence}, 2018.

\bibitem{bastani2021safe}
O.~Bastani.
\newblock Safe reinforcement learning with nonlinear dynamics via model predictive shielding.
\newblock {\em American Control Conference}, 2021.

\bibitem{konighofer2020shield}
B.~K\"onighofer, F.~Lorber, N.~Jansen, and R.~Bloem.
\newblock Shield synthesis for reinforcement learning.
\newblock {\em Leveraging Applications of Formal Methods (ISoLA)}, 2020.

\bibitem{sha2001using}
L.~Sha.
\newblock Using simplicity to control complexity.
\newblock {\em IEEE Software}, 18(4):20--28, 2001.

\bibitem{leucker2009brief}
M.~Leucker and C.~Schallhart.
\newblock A brief account of runtime verification.
\newblock {\em Journal of Logic and Algebraic Programming}, 78(5):293--303, 2009.

\bibitem{dibaji2019systems}
S.~M.~Dibaji, M.~Pirani, D.~B.~Flamholz, A.~M.~Annaswamy, K.~H.~Johansson, and A.~Chakrabortty.
\newblock A systems and control perspective of {CPS} security.
\newblock {\em Annual Reviews in Control}, 47:394--411, 2019.

\bibitem{chandola2009anomaly}
V.~Chandola, A.~Banerjee, and V.~Kumar.
\newblock Anomaly detection: A survey.
\newblock {\em ACM Computing Surveys}, 41(3):1--58, 2009.

\bibitem{chen2020review_bess}
T.~Chen, Y.~Jin, H.~Lv, A.~Yang, M.~Liu, B.~Chen, Y.~Xie, and Q.~Chen.
\newblock Applications of lithium-ion batteries in grid-scale energy storage systems.
\newblock {\em Transactions of Tianjin University}, 26:208--217, 2020.

\bibitem{pineau2021improving}
J.~Pineau, P.~Vincent-Lamarre, K.~Sinha, V.~Larivi\`ere, A.~Beygelzimer, F.~d'Alch\'e-Buc, E.~Fox, and H.~Larochelle.
\newblock Improving reproducibility in machine learning research.
\newblock {\em Journal of Machine Learning Research}, 22(164):1--20, 2021.

\end{thebibliography}

% ===== APPENDIX ==============================================================
\appendix

\section{Hyperparameters and Configuration}\label{app:hyperparams}

\PaperTableCSV{TBL06_HYPERPARAMS}{Core hyperparameters and configuration settings for reproducibility.}

Key configuration values:
\begin{itemize}[leftmargin=*,topsep=2pt,itemsep=1pt]
  \item Forecast: 24-hour horizon, 168-hour lookback, 10-fold CV with 24-hour gap.
  \item Conformal: $\alpha = 0.10$, volatility binning window = 168 hours.
  \item RAC-Cert: $\kappa_r = 0.5$, $\kappa_d = 0.3$, $\kappa_s = 0.2$, shrinkage $\tau = 0.95$.
  \item Drift detection: KS $p$-value threshold = 0.01, degradation threshold = 0.15.
  \item Optimization: cost weight $\lambda_c = 1.0$, carbon weight $\lambda_g = 1.2$, retraining cadence = 30 days.
\end{itemize}

\section{Additional Figures}\label{app:figures}

\PaperFigure{FIG09_REGION_DATASET_CARDS}{Dataset cards for all four regions in the release scope.}

\PaperFigure{FIG11_TRANSFER_GENERALIZATION}{Cross-region transfer and generalization performance summary.}

\section{Reproducibility Checklist}\label{app:reproducibility}

All results can be reproduced with the following commands:
\begin{verbatim}
# Train all models and build uncertainty artifacts
make train-all-ba

# Run CPSBench fault sweep
make r1-cpsbench

# Run ablation study
make r1-full

# Build publication artifacts and verify claims
make r1-verify
make paper-sync
\end{verbatim}

Canonical run IDs: DE \texttt{\PaperMetric{DE_STOCH_RUN_ID}{20260217\_165756}}, US \texttt{\PaperMetric{US_STOCH_RUN_ID}{20260217\_182305}}.  All artifacts are hash-tracked in the release manifest.

\section{Operational Failure Modes}\label{app:failure_modes}

\begin{table}[htbp]
\centering\small
\caption{Operational failure modes and safety controls.}
\label{tab:failure_modes}
\begin{tabular}{p{2.5cm}p{2.2cm}p{2.5cm}p{2.5cm}p{2.3cm}}
\toprule
Failure Mode & Trigger & Detection & Mitigation & Fallback \\
\midrule
Missing/stale data & Upstream delay & Readiness checks & Block unsafe dispatch & Hold current state \\
Forecast drift & Distribution shift & KS/model monitors & Retraining workflow & Conservative mode \\
Solver infeasible & Extreme intervals & Solver status flags & Safe infeasible response & Grid-only baseline \\
Unsafe command & SOC/power violation & BMS validation & Reject with detail & Maintain safe state \\
Control degradation & Heartbeat timeout & Watchdog detection & Lock remote control & Manual/local \\
\bottomrule
\end{tabular}
\end{table}

\end{document}
